\documentclass[10pt,a4paper]{amsart}
\usepackage{amsmath,amssymb,amsthm,graphicx}
% Retained minimal preamble for the Stacks Project chapter more-algebra.tex (More on Algebra).
% The chapter file pinned-original-source.tex carries only \input{preamble} on its first line, so no
% preamble text lives inside the chapter itself.  The Stacks Project uses one single shared preamble for
% every chapter; the snapshot stages it identically as sources/stacks-chow/preamble.tex and
% sources/stacks-etale-cohomology/preamble.tex.  The block below is copied mechanically, line for line,
% from sources/stacks-chow/preamble.tex lines 14-18 (the Xy-pic setup) and 154-217 (the native theorem
% declarations together with the mathematical macros).  Nothing was added, removed or redefined; the
% wrapper class and the amsmath/amssymb/amsthm/graphicx base packages are supplied by the assembler.
% Chapter: more-algebra.tex -- More on Algebra.
% For commutative diagrams we use Xy-pic
\usepackage[all]{xy}

% We use 2cell for 2-commutative diagrams.
\xyoption{2cell}

% Theorem environments.
%
\theoremstyle{plain}
\newtheorem{theorem}[subsection]{Theorem}
\newtheorem{proposition}[subsection]{Proposition}
\newtheorem{lemma}[subsection]{Lemma}

\theoremstyle{definition}
\newtheorem{definition}[subsection]{Definition}
\newtheorem{example}[subsection]{Example}
\newtheorem{exercise}[subsection]{Exercise}
\newtheorem{situation}[subsection]{Situation}

\theoremstyle{remark}
\newtheorem{remark}[subsection]{Remark}
\newtheorem{remarks}[subsection]{Remarks}

\numberwithin{equation}{subsection}

% Macros
%
\def\lim{\mathop{\mathrm{lim}}\nolimits}
\def\colim{\mathop{\mathrm{colim}}\nolimits}
\def\Spec{\mathop{\mathrm{Spec}}}
\def\Hom{\mathop{\mathrm{Hom}}\nolimits}
\def\Ext{\mathop{\mathrm{Ext}}\nolimits}
\def\SheafHom{\mathop{\mathcal{H}\!\mathit{om}}\nolimits}
\def\SheafExt{\mathop{\mathcal{E}\!\mathit{xt}}\nolimits}
\def\Sch{\mathit{Sch}}
\def\Mor{\mathop{\mathrm{Mor}}\nolimits}
\def\Ob{\mathop{\mathrm{Ob}}\nolimits}
\def\Sh{\mathop{\mathit{Sh}}\nolimits}
\def\NL{\mathop{N\!L}\nolimits}
\def\CH{\mathop{\mathrm{CH}}\nolimits}
\def\proetale{{pro\text{-}\acute{e}tale}}
\def\etale{{\acute{e}tale}}
\def\QCoh{\mathit{QCoh}}
\def\Ker{\mathop{\mathrm{Ker}}}
\def\Im{\mathop{\mathrm{Im}}}
\def\Coker{\mathop{\mathrm{Coker}}}
\def\Coim{\mathop{\mathrm{Coim}}}

% Boxtimes
%
\DeclareMathSymbol{\boxtimes}{\mathbin}{AMSa}{"02}

%
% Macros for moduli stacks/spaces
%
\def\QCohstack{\mathcal{QC}\!\mathit{oh}}
\def\Cohstack{\mathcal{C}\!\mathit{oh}}
\def\Spacesstack{\mathcal{S}\!\mathit{paces}}
\def\Quotfunctor{\mathrm{Quot}}
\def\Hilbfunctor{\mathrm{Hilb}}
\def\Curvesstack{\mathcal{C}\!\mathit{urves}}
\def\Polarizedstack{\mathcal{P}\!\mathit{olarized}}
\def\Complexesstack{\mathcal{C}\!\mathit{omplexes}}
% \Pic is the operator that assigns to X its picard group, usage \Pic(X)
% \Picardstack_{X/B} denotes the Picard stack of X over B
% \Picardfunctor_{X/B} denotes the Picard functor of X over B
\def\Pic{\mathop{\mathrm{Pic}}\nolimits}
\def\Picardstack{\mathcal{P}\!\mathit{ic}}
\def\Picardfunctor{\mathrm{Pic}}
\def\Deformationcategory{\mathcal{D}\!\mathit{ef}}
\title{A canonical isomorphism between the determinant of the middle term of a short exact sequence of finite projective modules and the tensor product of the outer determinants}
\author{The Stacks Project Authors}
\date{stacks-project:more-algebra.tex:local-snapshot-20261012; GNU FDL 1.2 or later}
\begin{document}
\maketitle
\setcounter{subsection}{1}
\noindent\textit{Editorial note.} All delivered mathematical text is the author's own, copied from the pinned Stacks Project chapter file more-algebra.tex; no statement or proof step was rewritten, reordered or removed.\par\medskip
\noindent\textit{Editorial note.} The chapter file pinned-original-source.tex carries only \textbackslash{}input\{preamble\} on its first line, so no preamble text lives inside the chapter. The retained file sources/stacks-more-algebra/retained-preamble.tex is a mechanical, line-for-line copy of the project's single shared preamble as staged in the snapshot; nothing was added, removed or redefined and no package was added.\par\medskip
\noindent\textit{Editorial note.} The source URL is the source-level Stacks Project URL, as required when more than one card is produced from one source; the per-claim tag is carried by claim\_key.\par\medskip
\noindent\textit{Editorial note.} The proof environment is opened by the author as \textbackslash{}begin\{proof\}[First proof] and the author writes 'We omit the detailed proof that this is well defined and an isomorphism'; the environment and that sentence are kept verbatim.\par\medskip
\noindent\textit{Allowed standard prerequisites (editorial).}\begin{itemize}
\item Determinant (top exterior power) of a finite projective module and local constancy of rank. the notation det in the statement at source lines 36291-36292 and the rank decomposition at proof lines 36296-36299. Supplies the two sides of the required isomorphism and the product decomposition that makes the ranks constant on each factor.
\item Exterior powers of modules and induced maps on top exterior powers. the map wedge\textasciicircum{}i(M') tensor wedge\textasciicircum{}j(M'') -> wedge\textasciicircum{}\{i+j\}(M) at proof lines 36301-36305. Is the explicit form the canonical isomorphism takes once the ranks are constant.
\end{itemize}\medskip
\begin{lemma}
\label{lemma-det-ses}
Let $R$ be a ring. Let
$$
0 \to M' \to M \to M'' \to 0
$$
be a short exact sequence of finite projective $R$-modules. Then there
is a canonical isomorphism
$$
\gamma : \det(M') \otimes \det(M'') \longrightarrow \det(M)
$$
\end{lemma}

\begin{proof}[First proof]
First proof. Decompose $R$ into a product of rings $R_{ij}$ such that
$M' = \prod M'_{ij}$ and $M'' = \prod M''_{ij}$ where
$M'_{ij}$ has rank $i$ and $M''_{ij}$ has rank $j$.
Of course then $M = \prod M_{ij}$ and $M_{ij}$ has rank $i + j$.
This reduces us to the case where $M'$ and $M''$ have constant
rank say $i$ and $j$. In this case we have to construct a canonical
map
$$
\wedge^i(M') \otimes \wedge^j(M'') \longrightarrow \wedge^{i + j}(M)
$$
To do this choose $m'_1, \ldots, m'_i$ in $M'$ and
$m''_1, \ldots, m''_j$ in $M''$. Denote $m_1, \ldots, m_i \in M$
the images of $m'_1, \ldots, m'_i$ and denote
$m_{i + 1}, \ldots , m_{i + j} \in M$ elements mapping to
$m''_1, \ldots, m''_j$ in $M''$. Our rule will be that
$$
m'_1 \wedge \ldots \wedge m'_i \otimes
m''_1 \wedge \ldots \wedge m''_j
\longmapsto
m_1 \wedge \ldots \wedge m_{i + j}
$$
We omit the detailed proof that this is well defined and an isomorphism.
\end{proof}

\end{document}
